\documentclass[a4paper,12pt]{article}
\usepackage[utf8]{inputenc}
\usepackage{geometry}
\usepackage{fancyhdr}
\usepackage{xcolor}
\usepackage{minted}
\usepackage{graphicx}
\usepackage{hyperref}
\usepackage{enumitem}

% Geometry setup
\geometry{top=2.5cm, bottom=2.5cm, left=2.5cm, right=2.5cm}

% Header and Footer
\pagestyle{fancy}
\fancyhf{}
\lhead{\textbf{CMP9134: Software Engineering}}
\rhead{\textbf{Week 2}}
\lfoot{University of Lincoln}
\rfoot{Page \thepage}

% Color definitions
\definecolor{lightgray}{rgb}{0.95,0.95,0.95}

\begin{document}

%----------------------------------------------------------------------------------------
%   TITLE SECTION
%----------------------------------------------------------------------------------------

\begin{center}
    \LARGE \textbf{Lab Sheet 2: Agile Working} \\
    \vspace{0.5cm}
    \large \textbf{Objectives}
\end{center}

\noindent By completing today's workshop, you will:
\begin{itemize}
    \item Configure a \textbf{Trello Board} specifically for your Assessment 1 (Robot System).
    \item Practice \textbf{"Personal Scrum"} by adapting Agile roles for a solo developer.
    \item Breakdown the Assessment Brief into \textbf{User Stories} and \textbf{Tasks}.
    \item Integrate \textbf{GitHub Issues} with Trello to link your planning with your code.
\end{itemize}

\vspace{0.5cm}
\hrule
\vspace{0.5cm}

\section*{Introduction: The Solo Agile Mindset}
While Scrum is typically a team sport, its principles are vital for managing individual coursework. In this module, you are wearing three hats simultaneously:
\begin{enumerate}
    \item \textbf{The Product Owner:} You decide \textit{what} to build (based on the Assessment Brief).
    \item \textbf{The Scrum Master:} You ensure you are working efficiently and removing distractions.
    \item \textbf{The Developer:} You write the code and tests.
\end{enumerate}
Today, we set up the tools to manage this workflow.

%----------------------------------------------------------------------------------------
%   TASK 1
%----------------------------------------------------------------------------------------

\section*{Task 1: Project Board Setup}

\textbf{Goal:} Create a Trello board tailored for the Robot Management System.

\begin{enumerate}
    \item \textbf{Log in to Trello:} Go to \href{https://trello.com}{trello.com} (use the account created in the lecture/workshop).
    \item \textbf{Create a New Board:}
    \begin{itemize}
        \item Click \textbf{Create new board}.
        \item Title: \texttt{CMP9134 - Robot Management System}.
        \item Visibility: Workspace or Public (Private is fine, but you may need want to make it public to add to your assessment report for the final submission).
    \end{itemize}
    \item \textbf{Configure Lists for Solo Work:}
    Rename the standard lists to match a "Personal Kanban" flow. Create the following columns left-to-right (Click on add list):
    \begin{itemize}
        \item \textbf{Product Backlog:} (All the ideas and requirements from the Brief).
        \item \textbf{Sprint Backlog (This Week):} (What you realistically plan to finish this week).
        \item \textbf{In Progress (Today):} (The 1 or 2 tasks you are working on \textit{right now}. Limit this to max 2!).
        \item \textbf{Waiting/Blocked:} (Tasks waiting for external help, e.g., "Ask lecturer about API").
        \item \textbf{Done:} (Completed tasks).
    \end{itemize}
\end{enumerate}

%----------------------------------------------------------------------------------------
%   TASK 2
%----------------------------------------------------------------------------------------

\section*{Task 2: From Assessment Brief to Backlog}

\textbf{Goal:} Breakdown the Assessment 1 requirements into Agile User Stories.

\noindent Open the \textbf{Assessment 01 Framework Brief} (available on Blackboard). You need to extract the requirements and turn them into cards in your \textbf{Product Backlog} list.

\vspace{0.3cm}
\noindent \textbf{Template for User Stories:}
\textit{"As a $<$Role$>$, I want $<$Feature$>$, so that $<$Benefit$>$."}

\vspace{0.3cm}
\noindent \textbf{Action:} Create at least \textbf{5 cards} in your Backlog based on the brief. Use the examples below as a guide:

\begin{enumerate}
    \item \textbf{Authentication (Security):}
    \textit{"As a User, I want to log in with a secure password, so that unauthorized people cannot control the robot."}
    \item \textbf{RBAC (Access Control):}
    \textit{"As a 'Viewer', I want to see the robot status but NOT move it, so that I don't accidentally disrupt a mission."}
    \item \textbf{Dashboard (UI):}
    \textit{"As a Commander, I want to see the robot's location on a 2D grid, so that I can plan its path."}
    \item \textbf{Telemetry (API Integration):}
    \textit{"As a User, I want to see the current battery level, so I know when the robot needs charging."}
    \item \textbf{Logging (Reliability):}
    \textit{"As an Auditor, I want a log of all commands sent, so that I can investigate accidents."}
\end{enumerate}

\noindent \textit{Tip: Don't forget non-functional requirements! Add cards for "Set up Docker Container", "Configure CI/CD Pipeline", and "Write AI Verification Log".}

%----------------------------------------------------------------------------------------
%   TASK 3
%----------------------------------------------------------------------------------------

\section*{Task 3: GitHub Integration (Issues \& Trello)}

\textbf{Goal:} Link your project management (Trello) with your code repository (GitHub).

\noindent \textbf{Why?} In professional environments, you never write code without a ticket. Linking them allows you to click a card in Trello and see the exact Pull Request or Branch associated with it.

\begin{enumerate}
    \item \textbf{Create your Assessment Repository if you haven't already in Week 1!}).
    \item \textbf{Enable the GitHub Power-Up in Trello:}
    \begin{itemize}
        \item On your Trello Board, click the \textbf{... Menu} (top right) $\rightarrow$ \textbf{Power-Ups} (the small shuttle icon).
        \item Search for \textbf{GitHub}.
        \item Click \textbf{Add} and then \textbf{Settings} to authorize your GitHub account.
        \item Click \textbf{Edit Power-Up Settings}, \textbf{Add Repo}, then link your GitHub account. 
        \item Select your \texttt{CMP9134\_Assessment1} repository to link it.
    \end{itemize}

    \item \textbf{The Workflow (to replicate for each new issue and card pair):}
    \begin{enumerate}
        \item Go to your GitHub repository and create a new \textbf{Issue} (e.g., add title "Implement Authentication", description and any other details). Issues are a way to plan, discuss, and track a specific part of your work. Use the following guidance for creating the issue: \url{https://docs.github.com/en/issues/tracking-your-work-with-issues/learning-about-issues/quickstart}.
        \item Fill in the description with the user story and acceptance criteria.
        \item Save the Issue.
        \item Open a card in Trello, or select one of the previously defined ones (e.g., the "Authentication" story).
        \item Click the \textbf{Power-Ups} button on the bottom side of the card.
        \item Click the \textbf{GitHub} button.
        \item Select \textbf{Attach Issue} and select the corresponding GitHub Issue you just created.
        \item You should now see the GitHub Issue linked to your Trello card. You can click on it to view details or navigate to GitHub directly.
    \end{enumerate}
    
    \item \textbf{Linking Branches (Optional but Recommended):}
    When you start coding, you can create a branch specifically for that issue (e.g., \texttt{feature/auth-login}). You can attach this branch to the Trello card using the same GitHub button, giving you full traceability.
\end{enumerate}

%----------------------------------------------------------------------------------------
%   TASK 4
%----------------------------------------------------------------------------------------

\section*{Task 4: The "Personal Scrum" Ritual}

\textbf{Goal:} Plan your first Sprint.

\noindent You don't have a team to hold meetings with, so you must hold them with yourself.

\begin{enumerate}
    \item \textbf{Sprint Planning:}
    Look at your \textbf{Product Backlog}. Move the top 3 items you want to achieve \textit{this week} into \textbf{Sprint Backlog}.
    \textit{(e.g., "Set up Docker Environment" and "Basic Hello World API Connection").}
    
    \item \textbf{Daily Stand-up (Morning Routine):}
    Every time you sit down to work on this module:
    \begin{itemize}
        \item Look at \textbf{In Progress}. Is it empty?
        \item If yes, pull \textbf{ONE} card from \textbf{Sprint Backlog} into \textbf{In Progress}.
        \item If no, finish that card before starting a new one!
    \end{itemize}
    
    \item \textbf{Definition of Done:}
    Before moving a card to \textbf{Done}, verify:
    \begin{itemize}
        \item [ ] Code is written.
        \item [ ] Code is pushed to GitHub.
        \item [ ] (If using AI) AI Verification Log is updated.
        \item [ ] It works locally.
    \end{itemize}
\end{enumerate}

%----------------------------------------------------------------------------------------
%   RESOURCES
%----------------------------------------------------------------------------------------

\section*{Further Resources}

\begin{itemize}
    \item \textbf{Writing Good User Stories:} \href{https://www.atlassian.com/agile/project-management/user-stories}{https://www.atlassian.com/agile/project-management/user-stories}
    \item \textbf{Trello GitHub Power-Up Guide:} \href{https://support.atlassian.com/trello/docs/using-the-github-power-up/}{https://support.atlassian.com/trello/docs/using-the-github-power-up/}
    % \item \textbf{Personal Kanban:} \href{https://personalkanban.com/}{https://personalkanban.com/} - A method for managing individual workload.
    \item \textbf{Assessment Brief:} Check Blackboard Assessment 1 folder for details regarding the Robot Management System.
\end{itemize}

\end{document}